\section{Algorithmic Checking: \smtlan}
\label{sec:theory:algorithmic}
\label{sec:algorithmic}

Next, we describe \smtlan, a conservative, first order
approximation of \corelan where higher order features are
approximated with uninterpreted functions
and the undecidable type subsumption rule \rsubbasecore
is replaced with a decidable one \rsubbasesmt, yielding
an SMT-based algorithmic type system that is both sound and
decidable.

\mypara{Syntax: Terms \& Sorts}
%
Figure~\ref{fig:smtsyntax} summarizes the syntax
of \smtlan, the \emph{sorted} (SMT-)
decidable logic of quantifier-free equality,
uninterpreted functions and linear
arithmetic (QF-EUFLIA) ~\citep{Nelson81,SMTLIB2}.
%
The \emph{terms} of \smtlan include
integers $n$,
booleans $b$,
variables $x$,
data constructors $\dc$ (encoded as constants),
fully applied unary \unop and binary \binop operators,
and application $x\ \overline{\pred}$ of an uninterpreted function $x$.
%
The \emph{sorts} of \smtlan include built-in
integer \tint and \tbool for representing
integers and booleans.
%
%% NV reflected functions and measures are first order
%% NV because
%% NV 1. they can be partially applied
%% NV 2. they can be passed as arguments
The interpreted functions of \smtlan, \eg
the logical constants $=$ and $<$,
%% NV and the uninterpreted functions app and lam
%% NV but we have not introduced these yet
have the function sort $\sort \rightarrow \sort$.
%
Other functional values in \corelan, \eg
reflected \corelan functions and
$\lambda$-expressions, are represented as
first-order values with
the uninterpreted sort \tsmtfun{\sort}{\sort}.
%
%%The uninterpreted functions of \smtlan, which
%%correspond to reflected \corelan functions,
%%have the function sort $\sort \rightarrow \sort$.
%%%
%%Other functional values in \corelan, \eg
%%$\lambda$-expressions, are represented as
%%first-order values in \smtlan with
%%uninterpreted sort \tsmtfun{\sort}{\sort}.
%%%
The universal sort \tuniv represents all other values.

\mypara{Semantics: Satisfaction \& Validity}
%
An assignment $\sigma$ maps
variables to terms
%
$\sigma \defeq \{ \assignto{x_1}{\pred_1}, \ldots, \assignto{x_n}{\pred_n}\}$.
%
We write
%
${\sigma \models \pred}$
%
if the assignment $\sigma$ is a
\emph{model of} $\pred$, intuitively
if $\applysub{\sigma}{\pred}$ ``is true''~\cite{Nelson81}.
%
A predicate $\pred$ \emph{is satisfiable} if
there exists ${\sigma\models\pred}$.
%
A predicate $\pred$ \emph{is valid} if
for all assignments ${\sigma\models\pred}$.


\subsection{Transforming \corelan into \smtlan}
%
\label{subsec:embedding}

The judgment
\tologicshort{\env}{e}{\typ}{\pred}{\sort}{\smtenv}{\axioms}
states that a $\corelan$ term $e$ is transformed,
under an environment $\env$, into a
$\smtlan$ term $\pred$.
%
Most of the transformation rules are identity
and can be found in~\cite{appendix}.
Here we discuss the non-identity ones.

\mypara{Embedding Types}
%
We embed \corelan types into \smtlan sorts as:
%
\[\arraycolsep=1.4pt
\begin{array}{rclcrcl}
\embed{\tint}                       & \defeq &  \tint  & \quad &
                                             &         &          \\
\embed{\tbool}                      & \defeq &  \tbool & \quad &
\embed{\tref{v}{\btyp^{[\tlabel]}}{\reft}} &\defeq & \embed{\btyp}\\
\embed{T}                           & \defeq &  \tuniv & \quad &
\embed{\tfun{x}{\typ_x}{\typ}} & \defeq & \tsmtfun{\embed{\typ_x}}{\embed{\typ}}\\
\end{array}
\]


\mypara{Embedding Constants}
%
Elements shared on both \corelan and \smtlan
translate to themselves.
%
These elements include
booleans,
integers,
variables,
binary
and unary
operators.
%
SMT solvers do not support currying,
and so in \smtlan, all function symbols
must be fully applied.
%
Thus, we assume that all applications
to primitive constants and data
constructors are fully applied, \eg by converting
source terms like @(+ 1)@ to
@(\z -> z + 1)@.
%


\mypara{Embedding Functions}
%
As \smtlan is first-order, we
embed $\lambda$-abstraction
using the uninterpreted function
\smtlamname{}{}.
%
$$
\inference{
    \tologicshort{\env, \tbind{x}{\typ_x}}{e}{}{\pred}{}{}{} &
        \hastype{\env}{(\efun{x}{}{e})}{(\tfun{x}{\typ_x}{\typ})}\\
}{
	\tologicshort{\env}{\efun{x}{}{e}}{(\tfun{x}{\typ_x}{\typ})}
	        {\smtlamname{\embed{\typ_x}}{\embed{\typ}}\ {x}\ {\pred}}
	        {\sort'}{\smtenv, \tbind{f}{\sort'}}{\andaxioms{\{\axioms_{f_1}, \axioms_{f_2}\}}{\axioms}}
}[\lgfun]
$$
%
The term $\efun{x}{}{e}$ of type
${\typ_x \rightarrow \typ}$ is transformed
to
${\smtlamname{\sort_x}{\sort}\ x\ \pred}$
of sort
${\tsmtfun{\sort_x}{\sort}}$, where
%
$\sort_x$ and $\sort$ are respectively
$\embed{\typ_x}$ and $\embed{\typ}$,
%
${\smtlamname{\sort_x}{\sort}}$
is a special uninterpreted function
of sort
${\sort_x \rightarrow \sort\rightarrow\tsmtfun{\sort_x}{\sort}}$,
and
$x$ of sort $\sort_x$ and $r$ of sort $\sort$ are
the embedding of the binder and body, respectively.
%
As $\smtlamname{}{}$ is an SMT-function,
it \emph{does not} create a binding for $x$.
%
Instead, $x$ is renamed to
a \emph{fresh} name pre-declared in
the SMT logic.


\mypara{Embedding Applications}
%
We embed applications via
defunctionalization~\citep{Reynolds72}
using the uninterpreted $\smtappname{}{}$:
%
$$
\inference{
	\tologicshort{\env}{e'}{\typ_x}{\pred'}{\embed{\typ_x}}{\smtenv}{\axioms'}
	&
	\tologicshort{\env}{e}{\tfun{x}{\typ_x}{\typ}}{\pred}{\tsmtfun{\embed{\typ_x}}{\embed{\typ}}}{\smtenv}{\axioms}
	&
	\hastype{\env}{e}{{\typ_x}\rightarrow{\typ}}
}{
	\tologicshort{\env}{e\ e'}{\typ}{\smtappname{\embed{\typ_x}}{\embed{\typ}}\ {\pred}\ {\pred'}}{\embed{\typ}}{\smtenv}{\andaxioms{\axioms}{\axioms'}}
}[\lgapp]
$$
%
The term ${e\ e'}$, where $e$ and $e'$ have
types ${\typ_x \rightarrow \typ}$ and $\typ_x$,
is transformed to
${\tbind{\smtappname{\sort_x}{\sort}\ \pred\ \pred'}{\sort}}$
where
%
$\sort$ and $\sort_x$ are $\embed{\typ}$ and $\embed{\typ_x}$,
the
${\smtappname{\sort_x}{\sort}}$
is a special uninterpreted function of sort
${\tsmtfun{\sort_x}{\sort} \rightarrow \sort_x \rightarrow \sort}$,
and
$\pred$ and $\pred'$ are the respective translations of $e$ and $e'$.


\mypara{Embedding Data Types}
%
We translate each data constructor to a
predefined \smtlan constant ${\smtvar{\dc}}$ of
sort ${\embed{\constty{\dc}}}$.
$$
	\tologicshort{\env}{\dc}{\constty{\dc}}{\smtvar{\dc}}{\embed{\constty{\dc}}}{\emptyset}{\emptyaxioms}
$$
%
For each datatype, we assume the existence of reflected functions that
\emph{check} the top-level constructor
and \emph{select} their individual fields.
%
For example, for lists, we assume the existence of measures:
%
\begin{mcode}
    isNil []     = True         isCons (x:xs) = True          sel1 (x:xs) = x
    isNil (x:xs) = False        isCons []     = False         sel2 (x:xs) = xs
\end{mcode}
%
Due to the simplicity of their syntax the above checkers and selectors
can be automatically instantiated in the logic
(\ie without actual calls to the reflected functions at source level)
using the measure mechanism of~\citet{Vazou15}.

To generalize, let $\dc_i$ be a data constructor such that
$$
\constty{\dc_i} \defeq \typ_{i,1} \rightarrow \dots \rightarrow \typ_{i,n} \rightarrow \typ
$$
Then the \emph{check function}
${\checkdc{{\dc_i}}}$ has the sort
$\tsmtfun{\embed{\typ}}{\tbool}$
and the \emph{select function}
${\selector{\dc}{i,j}}$ has the sort
$\tsmtfun{\embed{\typ}}{\embed{\typ_{i,j}}}$.
%

\mypara{Embedding Case Expressions}
%
We translate case-expressions
of \corelan into nested $\mathtt{if}$
terms in \smtlan, by using the check
functions in the guards and the
select functions for the binders
of each case.
$$
\inference{
	\tologicshort{\env}{e}{\typ_e}{\pred}{\embed{\typ_e}}{\smtenv}{\axioms} &&
	\tologicshort{\env}{e_i\subst{\overline{y_i}}{\overline{\selector{\dc_i}{}\ x}}\subst{x}{e}}{\typ}{\pred_i}{\embed{\typ}}{\smtenv}{\axioms_i}
}{
	\tologicshort{\env}{\ecase{x}{e}{\dc_i}{\overline{y_i}}{e_i}}{\typ}
	 {\eif{\smtappname{}{}\ \checkdc{\dc_1}\ \pred}{\pred_1}{\ldots} \ \mathtt{else}\ \pred_n}{\embed{\typ}}{\smtenv}
	 {\andaxioms{\axioms}{\axioms_i}}
}%[\lgcase]
$$
%
For example, the body of the list append function
%
\begin{code}
    []     ++ ys = ys
    (x:xs) ++ ys = x : (xs ++ ys)
\end{code}
%
is reflected into the \smtlan refinement:
%
$
\ite{\mathtt{isNil}\ \mathit{xs}}
    {\mathit{ys}}
    {\mathtt{sel1}\ \mathit{xs}\
       \dcons\
       (\mathtt{sel2}\ \mathit{xs} \ \mathtt{++}\  \mathit{ys})}.
$
%
We favor selectors to the axiomatic translation of
HALO~\citep{HALO} to avoid universally quantified
formulas and the resulting unpredictability.



\subsection{Correctness of Translation from \corelan to \smtlan}

Informally, the translation relation $\tologicshort{\env}{e}{}{\pred}{}{}{}$
is correct in the sense that if $e$ is a terminating boolean expression
then $e$ reduces to \etrue \textit{iff} $\pred$ is SMT-satisfiable
by a model that respects $\beta$-equivalence.

\begin{definition}[$\beta$-Model]\label{def:beta-model}
A $\beta-$model $\bmodel$ is an extension of a model $\sigma$
where $\smtlamname{}{}$ and $\smtappname{}{}$ satisfy the
axioms of $\beta$-equivalence:
%
$$
\forall\ x\ y\ e.\  \smtlamname{}{}\ x\ e \ =\ \smtlamname{}{}\ y\ (e\subst{x}{y})
\quad \quad \quad \quad
\forall\ x\ e_x\ e.\ \smtappname{}{}\ (\smtlamname{}{}\ x\ e)\ e_x \ =\   e\subst{x}{e_x}
$$
%
\end{definition}


\mypara{Semantics Preservation}
%
We define the translation of a \corelan term
into \smtlan under the empty environment as
${\embed{e} \defeq \pred}$
\textit{iff} ${\tologicshort{\emptyset}{\refa}{}{\pred}{}{}{}}$.
%
A \emph{lifted substitution}
$\theta^\perp$ is a set of models $\sigma$
where each ``bottom'' in the substitution
$\theta$ is mapped to an arbitrary logical
value of the respective sort~\citep{Vazou14}.
%
We connect the semantics of \corelan and
\smtlan via the following:
% terms can connect evaluation of boolean
% \corelan expression to \smtlan predicates.

\begin{theorem}\label{thm:embedding-general}
If ${\tologicshort{\env}{\refa}{}{\pred}{}{}{}}$,
then for every ${\sub\in\interp{\env}}$
and every ${\sigma\in {\sub^\perp}}$,
if $\evalsto{\applysub{\sub^\perp}{\refa}}{v}$
then $\sigma^\beta \models \pred = \embed{v}$.
\end{theorem}

% For Boolean expressions we specialize the above to

\begin{corollary}\label{thm:embedding}
If ${\hastype{\env}{\refa}{\tbool}}$, $e$ reduces to a value and
${\tologicshort{\env}{\refa}{\tbool}{\pred}{\tbool}{\smtenv}{\axioms}}$,
then for every ${\sub\in\interp{\env}}$
and every ${\sigma\in {\sub^\perp}}$,
$\evalsto{\applysub{\sub^\perp}{\refa}}{\etrue}$ iff
$\sigma^\beta \models \pred$.
\end{corollary}

\subsection{Decidable Type Checking}

We make the type checking from Figure~\ref{fig:typing} decidable
by checking subtyping via an SMT solver.

\mypara{Verification Conditions}
%
The implication or \emph{verification condition} (VC)
${\vcond{\env}{\pred}}$ is \emph{valid} only if the
set of values described by $\env$ is subsumed by
the set of values described by $\pred$.
%
$\env$ is embedded into logic by conjoining
(the embeddings of) the refinements of
provably terminating binders~\cite{Vazou14}:
%
$$
\embed{\env} \defeq \bigwedge_{x \in \env} \embed{\env, x} %\\
\quad \mbox{where we embed each binder as} \quad
\embed{\env, x} \defeq \begin{cases}
                           \pred  & \text{if } \env(x)=\tref{x}{\btyp^{\tlabel}}{e},\
                                    \tologicshort{\env}{e}{\btyp}{\pred}{\embed{\btyp}}{\smtenv}{\axioms} \\
                           \etrue & \text{otherwise}.
                         \end{cases}
$$
%
% VC checking, and thus type checking of \corelan, is decidable,
% as \smtlan is carefully restricted to SMT-decidable theories.

\mypara{Subtyping via Verification Conditions}
%
We make subtyping, and hence, typing decidable,
by replacing the denotational base subtyping
rule \rsubbasecore with the conservative,
algorithmic version \rsubbasesmt of Figure~\ref{fig:typing} that uses an SMT
solver to check the validity of the subtyping VC.
%
We use Corollary~\ref{thm:embedding} to prove
soundness of subtyping.
%
\begin{lemma}\label{lem:subtyping} %[Conservative Subtyping]
If {\aissubtype{\env}{\tref{v}{\btyp}{e_1}}{\tref{v}{\btyp}{e_2}}}
then {\issubtype{\env}{\tref{v}{\btyp}{e_1}}{\tref{v}{\btyp}{e_2}}}.
\end{lemma}

\mypara{Soundness of \smtlan}
%
We write $\ahastype{\env}{e}{\typ}$ for the judgments that
can be derived using the algorithmic subtyping rule
\rsubbasesmt instead of the denotational rule \rsubbasecore.
%
Lemma~\ref{lem:subtyping} implies the soundness of \smtlan.
%
\begin{theorem}[Soundness of \smtlan]\label{thm:soundness-smt}
If \ahastype{\env}{e}{\typ} then \hastype{\env}{e}{\typ}.
\end{theorem}
%
Note that we \emph{do not} require the $\beta$-equivalence axioms
for soundness: if a VC is valid \emph{without} the $\beta$-equivalence
axioms, \ie when $\smtlamname{}{}$ and $\smtappname{}{}$ are
uninterpreted, then the VC is trivially valid \emph{with}
the axioms, and hence, by Theorem~\ref{thm:soundness-smt}
the program is safe.

%% OK \NV{NEW}
%% But, we \emph{do} require the $\beta$-equivalence axioms
%% for precision: the typing rules of \smtlan would accept
%% more safe programs if the $\beta$-equivalence
%% axioms were assumed during validity checking.
%% %
%% Though assumption of the quantified $\beta$-equivalence axioms
%% would render validity checking undecidable.
%% %
%% In the next section we present an incomplete, yet sound and decidable
%% approach to increase precision of type checking by instantiation of
%% $\beta$-equivalence axioms.
%% OK \NV{END NEW}
